% Generated draft: requires mathematical and source review.
\documentclass[11pt]{article}
\usepackage[margin=25mm]{geometry}
\usepackage{fontspec}
\setmainfont{Latin Modern Roman}
\usepackage{amsmath,amssymb,mathtools,unicode-math}
\setmathfont{Latin Modern Math}
\usepackage{enumitem,needspace,xurl,hyperref}
\setlength{\parindent}{0pt}
\setlength{\parskip}{5pt}
\setlength{\emergencystretch}{3em}
\newcommand{\N}{\mathbb{N}}
\newcommand{\Z}{\mathbb{Z}}
\newcommand{\Q}{\mathbb{Q}}
\newcommand{\R}{\mathbb{R}}
\newcommand{\C}{\mathbb{C}}
\newcommand{\F}{\mathbb{F}}
\newcommand{\A}{\mathbb{A}}
\newcommand{\PP}{\mathbb{P}}
\newcommand{\E}{\mathbb{E}}
\newcommand{\eps}{\varepsilon}
\newcommand{\dd}{\,\mathrm{d}}
\begin{document}
\section*{UnsolvedMath ID: 30002348}
\textbf{Source problem code:} catalogue entry
\subsection{Definitions and mathematical statement}
The supplied problem asks whether, for random groups at density \(D<1/2\), the least genus of a homologically essential closed surface subgroup is \(O(n/\log n)\), where \(n\) is the relator length.

The following is a conditional precise formulation. Suppose that a random-group ensemble \((G_n)\) at a fixed admissible density \(D<1/2\) has been specified. Adopt the conventions that surfaces are connected, closed and orientable, have genus at least two, and that homological essentiality uses integral group homology. Let \(\Sigma_g\) denote such a surface of genus \(g\). An injective homomorphism \(\rho\colon\pi_1(\Sigma_g)\hookrightarrow G\) is homologically essential under these conventions if
\[
\rho_*[\Sigma_g]\ne 0\quad\text{in }H_2(G;\Z),
\]
where \([\Sigma_g]\in H_2(\pi_1(\Sigma_g);\Z)\) is the fundamental class under the identification of \(\Sigma_g\) with a classifying space for its fundamental group. Define
\[
m(G)=\min\bigl\{g\in\Z:g\ge2,\ \exists\rho\colon\pi_1(\Sigma_g)\hookrightarrow G\text{ with }\rho_*[\Sigma_g]\ne0\bigr\},
\]
with \(m(G)=+\infty\) if this set is empty.

Under these conventions, a possible precise interpretation of the question is: for every fixed admissible \(D<1/2\) and every fixed choice of the remaining ensemble parameters, does there exist a constant \(C>0\), independent of \(n\), such that
\[
\lim_{n\to\infty}\Pr\!\left(m(G_n)\le C\frac{n}{\log n}\right)=1?
\]
The limit is taken along the relator lengths allowed by the chosen ensemble; \(C\) may depend on \(D\) and its other fixed parameters.

This formulation requires confirmation: the supplied record does not specify the random-group model, the number of generators, the probabilistic quantifier, the meaning of ``simplest'', the orientability or permitted genera of the surfaces, or the coefficient system and target homology used for essentiality. The displayed formulation records one explicit interpretation, rather than resolving those omissions. The stated question requests an upper bound; the supplied description of this as an optimal order does not supply a matching lower-bound assertion.
\subsection{Short English statement}
For random groups with density \(D<1/2\), is the smallest genus of a homologically essential closed surface subgroup at most a constant times \(n/\log n\), as the relator length \(n\) grows? The ensemble, probability convention and precise meaning of homological essentiality require confirmation. The supplied record labels the problem open; that status has not been independently verified.
\subsection{Sources}
\begin{itemize}
\item Supplied local UnsolvedMath registry record, permanent ID \(30002348\), title ``Minimal-Genus Surface Subgroups of Random Groups'', Oberwolfach Reports Open Problems collection, exact source code \texttt{catalogue entry}. No individual proposer is identified; the supplied proposed year is 2013.
\item Supplied report citation: \emph{Geometric Structures in Group Theory} (2014), Oberwolfach Reports, pp.\ 1629--1675, DOI \texttt{10.4171/owr/2013/28}. \url{https://doi.org/10.4171/owr/2013/28}. The supplied title year is 2014, whereas the DOI and proposed-year metadata contain 2013; this discrepancy is unresolved.
\item Calegari--Walker, \emph{Random groups contain surface subgroups}. Supplied literature reference: \url{https://arxiv.org/abs/1304.2188}.
\item Calegari--Walker, \emph{Surface subgroups from linear programming}. Supplied literature reference: \url{https://arxiv.org/abs/1212.2618}.
\item Literature assessment dated exactly \texttt{2026-08-21}: the supplied curation labels the problem open and reports a linear-in-relator-length genus bound for essential surface subgroups, while stating that the conjectured optimal \(n/\log n\) order for the simplest essential surfaces has not been proved. This is supplied dated curation, not an independently verified literature assessment. Its model, probability and essentiality qualifications require checking against the cited works.
\item The supplied registry labels the statement \texttt{corrected\_verified} and reports a manual semantic, typographic and LaTeX review. These are supplied editorial claims; the report and linked papers have not been independently consulted for this draft.
\end{itemize}
\end{document}
